\section{Conclusion}\label{sec:conclusion}

In the model, the bidder pool in a takeover depends on the information that buyers have before they commit to
preparation. A
stronger incumbent moves the class value gap from the challenger's profit to target proceeds, and because a
liquidity order can match the investor's direction in the order slot, an investor who knows the challenger's
class keeps a residual advantage after every order. A larger target-payoff spread can make that advantage worth
the cost of trade, so the price can recruit a challenger that a weaker incumbent fails to attract.

The comparison of the weak, strong and very strong incumbent holds with unique outcomes on an open set of
parameters (Proposition~\ref{prop:ccc}). Under the default tie rules, a quiet market and an informative market
coexist exactly on a region of the target-payoff spread and the preparation threshold
(Proposition~\ref{prop:coexistence}), and a stronger incumbent can move the economy into and out of that region; in
the informative market, entry rises with the spread while the high class mixes between a purchase and abstention
and stays constant once it buys for sure. When information is costly, competition creates research: at an even
prior, when an indifferent investor may mix over its best orders and the preparation threshold lies above the prior,
research and entry rise together with the spread until research reaches its cap, in the equilibrium in which a
purchase after successful research recruits (Proposition~\ref{prop:research}), and the low class's selling does not
interrupt the rise.

Sale terms affect entry through the challenger's expected payment and through the information in the stock price.
Under the hypotheses of the main result, access to prices raises target proceeds and acquisition surplus net of
preparation costs against the strong incumbent, and, with a single value in each class and the incumbent's highest
value between them, in a bargaining family of payment rules with a zero reserve, the weight that the rule places
on the runner-up's value decides the direction in which competition moves the target-payoff spread. A reserve that
excludes the low class can make the price informative against both
incumbents; if informed trade is then full and only a purchase recruits at both strengths, entry is flat in
incumbent strength.

The next theoretical step is the seller's choice of terms with the whole pricing and trading continuation retained.
The empirical tests ask whether the challenger's preparation in visible sale processes responds to incumbent
strength when the price precedes the decision, and, where the investor always buys after good news and always sells
after bad news and the prior lies below the preparation threshold, whether a purchase raises the price by more than a sale lowers
it once each impact is weighted by how often its order occurs, as Section~\ref{sec:empirical} predicts where a
purchase recruits.
